% ============================================================
%  INTRODUCTION — IA, Deep Learning et Machine Learning
%  Julien Rolland — M2 Développement Fullstack
% ============================================================
\documentclass[aspectratio=169, 10pt]{beamer}
\input{../shared/preamble}

\title[IA, Deep Learning et Machine Learning]{IA, Deep Learning\\ et Machine Learning}
\subtitle{Introduction --- Roadmap \& Modalités d'évaluation}
\date{}

% ============================================================
\begin{document}
% ============================================================

\begin{frame}
  \titlepage
\end{frame}

% ============================================================
\section{Le module}
% ============================================================

\begin{frame}{Vue d'ensemble}
  \begin{columns}[T]
    \begin{column}{0.52\textwidth}
      \begin{block}{Objectifs}
        À la fin de ce module, vous serez capables de :
        \begin{itemize}
          \item Comprendre les \textbf{fondements mathématiques}
                du Machine Learning
          \item Implémenter des modèles \textbf{from scratch}
                (régression, perceptron, autograd)
          \item Utiliser \textbf{PyTorch} et \textbf{scikit-learn}
                pour des projets réels
          \item Déployer un modèle en \textbf{production}
                (FastAPI, Docker, MLOps)
          \item Travailler en autonomie sur un
                \textbf{challenge compétitif}
        \end{itemize}
      \end{block}
    \end{column}
    \begin{column}{0.44\textwidth}
      \begin{exampleblock}{Format}
        \begin{itemize}
          \item \textbf{6 jours} --- 42 heures
          \item \textbf{50\,\%} Théorie (slides)
          \item \textbf{50\,\%} Pratique (notebooks)
        \end{itemize}
      \end{exampleblock}
      \smallskip
      \begin{block}{Prérequis}
        \begin{itemize}
          \item Python intermédiaire
          \item Algèbre linéaire de base
                (vecteurs, matrices)
          \item Calcul différentiel (dérivées)
        \end{itemize}
      \end{block}
    \end{column}
  \end{columns}
\end{frame}

% ============================================================
\section{Roadmap}
% ============================================================

\begin{frame}{Roadmap --- 6 Jours}
  \begin{center}
  \begin{tikzpicture}[font=\tiny,
    am/.style={draw, rounded corners=3pt, align=center,
               minimum width=2.05cm, minimum height=0.62cm,
               fill=jedy_blue!15, line width=0.5pt},
    pm/.style={draw, rounded corners=3pt, align=center,
               minimum width=2.05cm, minimum height=0.62cm,
               fill=jedy_mid!15, line width=0.5pt},
    lbl/.style={font=\bfseries\scriptsize, color=jedy_blue}
  ]
    % Positions X (espacement 2.3cm, noeuds 2.05cm => total 13.1cm)
    \def\xa{0}   \def\xb{2.3} \def\xc{4.6}
    \def\xd{6.9} \def\xe{9.2} \def\xf{11.5}

    % Headers
    \foreach \x/\d in {\xa/J1, \xb/J2, \xc/J3, \xd/J4, \xe/J5, \xf/J6}{
      \node[lbl] at (\x, 1.35) {\d};
    }

    % AM blocks
    \node[am] at (\xa, 0.72) {Tooling\\Tenseurs};
    \node[am] at (\xb, 0.72) {Perceptron\\Logistique};
    \node[am] at (\xc, 0.72) {Autograd\\from Scratch};
    \node[am] at (\xd, 0.72) {Embeddings\\\& Attention};
    \node[am] at (\xe, 0.72) {MLOps\\\& Production};
    \node[am, fill=jedy_alert!20] at (\xf, 0.72) {QCM Final\\(80\,\%)};

    % PM blocks
    \node[pm] at (\xa, 0.08) {Régression\\Linéaire};
    \node[pm] at (\xb, 0.08) {Généralisation\\Scikit-Learn};
    \node[pm] at (\xc, 0.08) {PyTorch\\CNN};
    \node[pm] at (\xd, 0.08) {Transformers\\Fine-tuning};
    \node[pm] at (\xe, 0.08) {Challenge\\Kaggle};
    \node[pm, fill=jedy_example!30] at (\xf, 0.08) {Soutenances\\(20\,\%)};

    % Separator line
    \draw[dashed, color=gray!50, line width=0.4pt]
      (-1.1, 0.40) -- (12.6, 0.40);

    % Legend
    \node[font=\tiny, color=jedy_blue] at (-0.75, 0.72) {AM};
    \node[font=\tiny, color=jedy_mid]  at (-0.75, 0.08) {PM};
  \end{tikzpicture}
  \end{center}

  \vspace{6pt}
  \begin{columns}[T]
    \begin{column}{0.30\textwidth}
      \begin{block}{\small J1--J2 : Fondations}
        \scriptsize Maths, NumPy, régression, classification,
        biais-variance, cross-validation.
      \end{block}
    \end{column}
    \begin{column}{0.30\textwidth}
      \begin{block}{\small J3--J4 : Deep Learning}
        \scriptsize Rétropropagation, PyTorch, CNN, NLP,
        Transformers, fine-tuning.
      \end{block}
    \end{column}
    \begin{column}{0.30\textwidth}
      \begin{block}{\small J5--J6 : Production \& Éval}
        \scriptsize MLOps, Docker, FastAPI, challenge
        compétitif, soutenance, QCM.
      \end{block}
    \end{column}
  \end{columns}
\end{frame}

% ============================================================
\section{Évaluation}
% ============================================================

\begin{frame}{Modalités d'Évaluation}
  \begin{columns}[T]
    \begin{column}{0.52\textwidth}
      \begin{block}{QCM Final --- 80\,\% de la note}
        \begin{itemize}
          \item \textbf{Date :} J6 matin --- 23 mars 2026
          \item \textbf{Format :} Questions à choix multiples
          \item \textbf{Périmètre :} l'ensemble des concepts
                vus durant la semaine
          \item \textbf{Durée :} à préciser
        \end{itemize}
      \end{block}
      \smallskip
      \begin{alertblock}{Conseil}
        Les notebooks sont la meilleure préparation.
        Comprendre \textbf{pourquoi} le code fonctionne
        vaut plus que mémoriser des formules.
      \end{alertblock}
    \end{column}
    \begin{column}{0.44\textwidth}
      \begin{block}{Challenge Kaggle --- 20\,\% de la note}
        \begin{tabular}{rl}
          \textbf{10 pts} & Performance (leaderboard privé)\\
          \textbf{10 pts} & Soutenance flash J6 (5~min)\\
        \end{tabular}
        \vspace{4pt}
        \begin{itemize}
          \item Travail en \textbf{groupes}
          \item Condition minimale : battre la baseline
        \end{itemize}
      \end{block}
    \end{column}
  \end{columns}
\end{frame}

\begin{frame}{Supports de Cours}
  \begin{center}
    \begin{exampleblock}{\centering Tous les supports sont disponibles sur GitHub}
      \centering\vspace{4pt}
      \Large\ttfamily https://github.com/kauzarc/efrei-ia-ml-2026-03
      \vspace{4pt}
    \end{exampleblock}
  \end{center}
  \vspace{4pt}
  \begin{columns}[T]
    \begin{column}{0.48\textwidth}
      \begin{block}{Contenu du dépôt}
        \begin{itemize}
          \item Slides PDF de chaque session
          \item Notebooks corrigés
          \item Notebooks exercices (squelettes)
        \end{itemize}
      \end{block}
    \end{column}
    \begin{column}{0.48\textwidth}
      \begin{alertblock}{Recommandation}
        Travaillez directement dans les notebooks
        exercices --- les corrigés sont là en cas
        de blocage, pas en remplacement.
      \end{alertblock}
    \end{column}
  \end{columns}
\end{frame}

\end{document}
